% Copyright (c) 2026 Geovana Neves. Licensed under CC BY 4.0 (https://creativecommons.org/licenses/by/4.0/). See LICENSE-AND-AUTHORSHIP.md.
%
% 00-fts-overview/text/04-pyfs.tex
%
% Part 4: how pyflightstream carries the workflow of Part 3 in a workspace:
% the matrix, the three input files a row cites, the four stages of a run,
% the four workflows, and storage and sync. Every statement here is defined
% in Guides 1 to 7 and in the package pages [1-4]; this part only places it.

\section{pyflightstream on the workflow}

\begin{frame}{What pyflightstream is}
\begin{columns}[T]
\begin{column}{0.52\textwidth}
\begin{itemize}\small
  \item A Python driver for FlightStream, MIT licensed, installed from PyPI
        (\code{pip install pyflightstream}). FlightStream itself, and its
        licence, are separate.
  \item It \textbf{writes} the solver script from a command database per
        build, \textbf{runs} it, \textbf{reads} what the solver exported, and
        \textbf{post-processes} it into tables with their provenance.
  \item A study is a \textbf{workspace}: a folder with a run matrix and the
        input files the matrix cites. No Python is written to run it.
\end{itemize}
\end{column}
\begin{column}{0.44\textwidth}
\slidefig[0.56]{g00-pyfs-layers}
\end{column}
\end{columns}
\begin{takeaway}A row names things, and the package makes the link: a row never
carries a path.\end{takeaway}
\end{frame}

\begin{frame}{The workspace}
\slidefig[0.80]{guide-workspace-tree}
\end{frame}

\begin{frame}{The ten stages, and where each one lives}
\relax
{\ftstablesize\renewcommand{\arraystretch}{1.08}
\begin{tabularx}{\linewidth}{@{}P{0.25\linewidth}RP{0.10\linewidth}@{}}
\toprule
\textbf{Stage} & \textbf{In the workspace} & \textbf{Guide} \\
\midrule
1, 2 Scope and matrix & one row of \code{matriz.fs} per study point or sweep:
  \code{DESCRIPTION}, \code{FLIGHT\_CONDITION}, \code{SWEEP\_VALUES},
  \code{RUN} & 1 \\
3, 4 Geometry and mesh & made outside the package; the file sits in
  \code{inputs/geometries/}, the row names it in \code{GEOMETRY};
  \pyfs{pyfs-matrix inventory} reads its boundaries & 1, 2 \\
5 Boundary conditions & the trailing edge by file or by detection, base
  regions and wake nodes as keys; rotors and actuator discs in the
  reference file & 2, 3 \\
6 Solver setup & \code{inputs/setups/s<id>.toml} (column \code{SET}); lengths,
  moment point and frames in \code{inputs/references/r<id>.toml} (\code{REF});
  \code{SYMMETRY} and \code{WORKFLOW} on the row & 3, 4 \\
7 Execution & \pyfs{plan}, \pyfs{run}, \pyfs{collect}; \code{FS\_BUILD},
  \code{NCPUS}, \code{WALLTIME} on the row & 1, 7 \\
8 Output definition & \code{inputs/pproc/p<id>.toml} (\code{PPROC}): what the
  run exports, declared before it runs & 4, 5 \\
9, 10 Verification and analyses & \pyfs{post} writes \code{post/}: polars,
  series, sections, probes, provenance; FSI couples a structure & 5, 6 \\
\bottomrule
\end{tabularx}}
\end{frame}

\begin{frame}{The row names things, the files define them}
\slidefig[0.60]{guide-inputs-into-run}
\begin{takeaway}\code{REF}, \code{SET} and \code{PPROC} carry a file stem and
\code{GEOMETRY} names the mesh: stated once, in one file, and every row citing
it demonstrably used it.\end{takeaway}
\end{frame}

\begin{frame}{Reference, setup, pproc: three files a row cites}
\begin{columns}[T]
\begin{column}{0.32\textwidth}
\begin{block}{\code{r<id>.toml}}
\small \textbf{What the coefficients divide by, and about which point}:
reference area and lengths, the moment point, aliases for groups of
surfaces, custom frames, and one block per rotor or actuator disc
(Guide 3).
\end{block}
\end{column}
\begin{column}{0.32\textwidth}
\begin{block}{\code{s<id>.toml}}
\small \textbf{How hard to solve and when to stop}: iterations and
convergence, the boundary layer and its coupling, the model and its
initialisation, wake settings, and a raw table for any command without a key
(Guide 4).
\end{block}
\end{column}
\begin{column}{0.32\textwidth}
\begin{block}{\code{p<id>.toml}}
\small \textbf{What the run exports, and how post reads it}: loads, sections,
probes, plots, the log last; the groups and derived columns post reduces
with. The averaging window belongs to the row (Guides 4, 5).
\end{block}
\end{column}
\end{columns}
\vspace{2mm}
\begin{itemize}\small
  \item The flight condition is a cell of the row, not a file: Mach, speed,
        Reynolds number, altitude, a temperature offset, each a constraint on
        one flow state [1].
  \item Any key only some rows need goes in the row's free cell,
        \code{VAR\_NAMES\_VALUES}.
\end{itemize}
\end{frame}

\begin{frame}[fragile]{Four stages of a run}
\slidefig[0.40]{guide-stage-flow}
\begin{lstlisting}[style=ftsbase]
pyfs-matrix plan    matriz.fs --workspace . --fs-version 26.123   # no seat, no queue
pyfs-matrix run     matriz.fs --workspace . --fs-version 26.123   # solves, or submits
pyfs-matrix collect --workspace .                                 # cluster: settled outputs
pyfs-matrix post    matriz.fs --workspace .                       # rebuilds the products
\end{lstlisting}
{\small \pyfs{plan} validates every row against the build's command database
and writes the receipt \pyfs{run} refuses to start without (Guide 1).\par}
\end{frame}

\begin{frame}{Four workflows, one per kind of physics}
\relax
{\ftstablesize\renewcommand{\arraystretch}{1.10}
\begin{tabularx}{\linewidth}{@{}P{0.17\linewidth}P{0.27\linewidth}RP{0.12\linewidth}@{}}
\toprule
\textbf{\code{WORKFLOW}} & \textbf{What is solved} & \textbf{The physics of Part 2} & \textbf{Cost} \\
\midrule
\code{steady} & one converged steady solve per point; a sweep is one job &
  the steady problem: tangency, Kutta, the fixed wake & seconds per point \\
\code{unsteady} & a march in physical time, nothing turning & Kelvin and
  the starting vortex; the step as a wake dimension; a window to average &
  minutes \\
\code{unsteady\_rotor} & the blades turn, the wake is shed in time & the
  helical wake; degrees per step, revolutions, a per-revolution window; the
  reference for any rotor & minutes to hours \\
\code{qsteady\_rotor} & blades held still, the free stream turning about the
  shaft; sector or wheel & an isolated, axisymmetric rotor where the reduced
  frequency is small & seconds \\
\bottomrule
\end{tabularx}}
\vspace{1mm}
\begin{itemize}\small
  \item The row names its workflow; a key another workflow owns is refused by
        name at \pyfs{plan}. Costs are orders of magnitude on one
        workstation [20].
  \item Guide 1 carries the decision chart, the quasi-steady rotor in depth,
        and when its answer stands for the turning blade.
\end{itemize}
\end{frame}

\begin{frame}{Storage and sync}
\begin{columns}[T]
\begin{column}{0.48\textwidth}
\begin{block}{What a run leaves}
\begin{itemize}\small
  \item \code{sims/}: one folder per matrix row, \code{sim\_<POL>}, with the
        scripts and one datapoint folder per point: the saved simulation and
        the solver's exports.
  \item \code{post/}: the products, rebuilt by \pyfs{post} with no seat.
  \item \code{archive/}: whatever a rebuild replaced.
  \item \code{runs.json}: the record of every run, with its state.
\end{itemize}
\end{block}
\end{column}
\begin{column}{0.48\textwidth}
\begin{block}{Two workspaces, one main}
\begin{itemize}\small
  \item A workstation workspace for planning, one point, looking at results;
        a cluster workspace for long marches and many points.
  \item \pyfs{pyfs-matrix sync}, run on the main, brings the cluster's runs
        home at a chosen level; it previews unless told \code{--apply}, and
        never copies a mesh [2].
\end{itemize}
\end{block}
\end{column}
\end{columns}
\begin{takeaway}You never edit anything under \code{sims/} or \code{post/}. If a
number there is wrong, an input is wrong.\end{takeaway}
\end{frame}
